\section*{An author's retrospective: why these shapes mattered}
\phantomsection\label{sec:retrospective}
\textit{This opening follows Hilmir's retrospective account approved on 19 October 2026. It presents his recollection in editorial prose, not a transcript or an independently dated reconstruction. The mathematical commentary that follows retains the established definitions and results of edition v0.1.}~\cite{History26}

\subsection*{Conversations, frustration and a visual association}
Hilmir remembers enjoying conversations with GPT about the universe. At one point he asked the model what \emph{it} would like to investigate scientifically. GPT suggested quantum physics and referred him to a university lecture. Some of the explanations irritated him rather than helping him understand. In his account, that reaction became part of the reason to keep asking questions and to look for a way of picturing the subject that made sense to him. It describes his experience of those explanations; it is not a judgment about academic teaching as a whole.

A separate exploration began with symbolism in \emph{Cyberpunk 2077}. He read Greek mythology and discussed its symbolic relationships to quantum ideas with GPT. These interests met in a particularly immediate visual moment: when he saw an image of a hydrogen atom, he associated its appearance with the Flower of Life and the vesica piscis. The resemblance prompted a question about geometric constants: which constants occur naturally in which forms, and what might those appearances help him understand?

The retrospective does not identify the exact lecture or hydrogen image. Its value is the remembered route into the work: pleasure in conversation, difficulty with an explanation, and a shape that suggested another question. The symbolic association was a starting point for investigation. It did not by itself identify an atomic wavefunction with a two-circle construction or establish a quantum result.

\subsection*{A proposed order of forms}
His tentative hierarchy began with a point or zero. The circle was the first continuous enclosed form in his proposed ordering. Next came what he understood as two lines: edge-on views of two circles, with a vesica becoming visible when the viewing angle changed by $90^\circ$. He then considered the triangle, square, pentagon, hexagon, heptagon and octagon. This was his way of organizing the construction he was imagining, rather than a mathematical classification of all possible forms.

The triangle carried several connected meanings for him: movement, a Z-direction, and a possible picture of momentum and energy. He associated it with a $\pm45^\circ$ range and an above/below relation. Two triangles could combine into a skewed box within his vesica picture. These descriptions preserve what he was trying to see. They do not supply a motion equation, a chosen axis or a measured momentum-energy relation.

At that stage he treated the octagon as the final necessary operator in his proposed geometry of unity. ``Operator'' here records a role in his organizing picture; it does not declare a linear or recursive operator without its own definition. Nor did he mean that polygons end at eight sides. Much later he contemplated the decagon as an outer hull or shell, not as an additional fundamental operator. The nonagon's role remained undecided. Keeping that undecided place matters: the retrospective is not a completed hierarchy reconstructed after the fact.

He also recalls possible links between his counting convention and the number nine in Three-Way work. Those links remain tentative recollections, not established derivations. In a separate symbolic association, he related $144{,}000$ to $360\times400$. The arithmetic is exact; the symbolic significance he attached to it does not establish a gradient law. These numerical associations belonged to the questions he was pursuing, with different evidential status from the exact geometry described later.

\subsection*{Completing the overlap, then changing perspective}
In the mental construction described here, the two-circle vesica overlap comes first. Only after that overlap is complete does one circle become the reference from which he reasons beyond the centre and into three dimensions. That change of reference helped motivate the Tri-Octagon. It does not introduce a literal third parent circle into the older two-circle formulas.

He also pictured moving a right circle by $45^\circ$ to the left and a bottom circle by $45^\circ$ to the right, mirroring or flipping views, and changing perspective by $90^\circ$. These are qualitative descriptions of the image he was manipulating. No rotation axes, active Euclidean embedding or transformation from those movements to the welded module were specified in this account. The circle used as a later reference and the two circles that first define the overlap have different roles in the description.

His earlier quantum and dark-matter questions, Vesica-delay ideas, de Broglie comparisons and later SRG investigations were different views and questions pursued from this common geometric intuition. Their common origin does not prove a physical equivalence among them. Likewise, the historical suggestion that a changed perspective might expose three wavelength responses remains an association he proposed. It is not a native result of the Tri-Octagon geometry or a theorem that changing one's view multiplies waves.

\Needspace{8\baselineskip}
\subsection*{The roles assigned to three constants}
Within the same imagined construction, Hilmir assigned distinct roles to $\pi$, Euler's $e$ and the golden ratio $\ph$. The circle was associated with $\pi$. He pictured $e$ as a distance from zero or the origin to the circle. He associated $\ph$ with a spiral progressing inward and outward through a constructed cylinder. These roles describe the meaning he was trying to give the constants, rather than an ordering of their numerical sizes.

He wrote $V=\pi e\ph$ as a proposed structural or volume interpretation of that picture. In this historical usage, the letter $V$ records his intended interpretation; it does not turn the product into a proved Euclidean cylinder-volume formula. The account supplies no calibrated spatial unit or conversion to a total-universe size. In particular, replacing his remembered roles with a newly chosen radius and height would change the account. The product itself also occurs in the documented constant chain retained below, where the defined target $T$ and the historical coupling proposal have their own explicit mathematical meanings. Those later definitions do not retrospectively prove the original volume interpretation.

The interest was therefore not simply in obtaining a familiar number. He was trying to understand how a circle, a distance from an origin, and an inward/outward spiral might belong to one construction. This edition preserves that question alongside the distinctions needed to read the surviving equations accurately.

\subsection*{The pleasure of a precise half-$\pi$ result}
One result gave the geometric picture a particularly concrete point of contact. In \emph{The Vesica Delay Method}, the parent-circle radius is $r=1/2$, the separation is $d=2r\cos\theta$, and the overlap perimeter is $C_{\rm overlap}=4r\theta$. At $\theta=45^\circ=\pi/4$, the perimeter is exactly $\pi/2$. This is the specific half-$\pi$ relationship Hilmir recalls with delight.~\cite{Delay,History26}

Here the quantity is the length of the two arcs bounding the overlap in the stated radius normalization. It is not an angular-sector fraction, nor does it identify every appearance of $1/2$ in another equation. Section~2 retains the circle construction, its limiting configuration and the distinction between the lens and a separate imagined eye-circle. It also retains the source criticism of the earlier VPQW documents: their printed $C=\pi$ and $r=1/(2\pi)$ are inconsistent under the ordinary circumference formula. The later, consistent Delay Method convention does not silently correct those older statements.

\subsection*{Questions carried through a collaboration}
Hilmir pursued these ideas with GPT and later Codex, asking the models to undertake substantial mathematical, writing and computational tasks when appropriate. The questions, visual associations and choices of what to pursue were meaningful parts of his contribution. AI assistance helped develop and express the work, while the surviving equations and computations have to be assessed on their own stated assumptions and evidence.

The record does not support assigning an individual author to every equation in the old conversations, and model assistance or internal model review is not external peer review. The purpose of this retrospective is to make the author's thinking available to future readers. It does not certify the earlier quantum or dark-matter hypotheses as experimental truths, and it does not require every later equation to be reinterpreted through an initial image.

\section*{Reading the historical picture against the mathematics}
The following eight numbered sections retain the earlier companion's mathematical account, including its normalization, exact statements, constants evidence and scoped negative results. Their equation numbering is unchanged. The first section records the earlier retrospective of 17 October; the new opening above supplies the expanded account approved on 19 October. Recollection, a printed historical proposal, a geometric theorem and a bounded numerical comparison remain different kinds of evidence.

The later discussion of a decagonal enclosure offers a precise example of this distinction: the accepted spherical-edge synthesis treats a decagon as an optional outer comparison with its own orientation choice. That documented construction does not turn the author's earlier hierarchy into a theorem about necessary operators.~\cite[Sec.~14; App.~D]{SP} Similarly, the target $T=\pi e\ph$ is preserved with its source definitions and qualifications. The original two figures illustrate specified circle and shell geometry; they are not reconstructions of private sketches or illustrations of a newly asserted physical law.

